% Options for packages loaded elsewhere
\PassOptionsToPackage{unicode}{hyperref}
\PassOptionsToPackage{hyphens}{url}
\documentclass[
  ignorenonframetext,
]{beamer}
\newif\ifbibliography
\usepackage{pgfpages}
\setbeamertemplate{caption}[numbered]
\setbeamertemplate{caption label separator}{: }
\setbeamercolor{caption name}{fg=normal text.fg}
\beamertemplatenavigationsymbolsempty
% remove section numbering
\setbeamertemplate{part page}{
  \centering
  \begin{beamercolorbox}[sep=16pt,center]{part title}
    \usebeamerfont{part title}\insertpart\par
  \end{beamercolorbox}
}
\setbeamertemplate{section page}{
  \centering
  \begin{beamercolorbox}[sep=12pt,center]{section title}
    \usebeamerfont{section title}\insertsection\par
  \end{beamercolorbox}
}
\setbeamertemplate{subsection page}{
  \centering
  \begin{beamercolorbox}[sep=8pt,center]{subsection title}
    \usebeamerfont{subsection title}\insertsubsection\par
  \end{beamercolorbox}
}
% Prevent slide breaks in the middle of a paragraph
\widowpenalties 1 10000
\raggedbottom
\AtBeginPart{
  \frame{\partpage}
}
\AtBeginSection{
  \ifbibliography
  \else
    \frame{\sectionpage}
  \fi
}
\AtBeginSubsection{
  \frame{\subsectionpage}
}
\usepackage{iftex}
\ifPDFTeX
  \usepackage[T1]{fontenc}
  \usepackage[utf8]{inputenc}
  \usepackage{textcomp} % provide euro and other symbols
\else % if luatex or xetex
  \usepackage{unicode-math} % this also loads fontspec
  \defaultfontfeatures{Scale=MatchLowercase}
  \defaultfontfeatures[\rmfamily]{Ligatures=TeX,Scale=1}
\fi
\usepackage{lmodern}
\ifPDFTeX\else
  % xetex/luatex font selection
\fi
% Use upquote if available, for straight quotes in verbatim environments
\IfFileExists{upquote.sty}{\usepackage{upquote}}{}
\IfFileExists{microtype.sty}{% use microtype if available
  \usepackage[]{microtype}
  \UseMicrotypeSet[protrusion]{basicmath} % disable protrusion for tt fonts
}{}
\makeatletter
\@ifundefined{KOMAClassName}{% if non-KOMA class
  \IfFileExists{parskip.sty}{%
    \usepackage{parskip}
  }{% else
    \setlength{\parindent}{0pt}
    \setlength{\parskip}{6pt plus 2pt minus 1pt}}
}{% if KOMA class
  \KOMAoptions{parskip=half}}
\makeatother
\setlength{\emergencystretch}{3em} % prevent overfull lines
\providecommand{\tightlist}{%
  \setlength{\itemsep}{0pt}\setlength{\parskip}{0pt}}
\usepackage{bookmark}
\IfFileExists{xurl.sty}{\usepackage{xurl}}{} % add URL line breaks if available
\urlstyle{same}
\hypersetup{
  hidelinks,
  pdfcreator={LaTeX via pandoc}}

\author{\texorpdfstring{}{}}
\date{}

\begin{document}

\section{Conclusion}\label{conclusion}

\begin{frame}[allowframebreaks]{Summary of Contributions}
\protect\phantomsection\label{summary-of-contributions}
This paper has applied the Cognitive Integrity Framework (CIF)
\cite{friedman2026cogsec1} across ten critical domains, demonstrating
that Goal Hijacking is not a narrow linguistic exploit but a structural
corruption of autonomous decision-making. The specific contributions
are:

\textbf{C1: CIF-AD-OODA Integration Model.} We formalized the
integration of three complementary frameworks---CIF (defense
mechanisms), Axiomatic Design (structural analysis)
\cite{suh2001axiomatic}, and the OODA Loop (temporal dynamics)
\cite{boyd1987patterns}---into a unified analytical model for Goal
Hijacking. This model enables systematic domain analysis through a
standardized five-step template.

\textbf{C2: Universal Attack Pattern Taxonomy.} Through cross-domain
analysis, we identified three universal attack patterns---FR Polarity
Inversion, Constraint Relaxation, and Context Boundary Violation---that
characterize all Goal Hijacking attacks as specific manipulations of the
Axiomatic Design Matrix. FR Polarity Inversion is the most prevalent
(5/10 domains), revealing that the most effective attacks \emph{co-opt}
rather than \emph{disable} agent capabilities.

\textbf{C3: CIF Mechanism Coverage Validation.} We demonstrated that all
five canonical CIF mechanisms appear across the ten-domain portfolio,
with each mechanism serving as a primary defense in at least three
domains. No domain requires mechanisms beyond the CIF vocabulary, and no
single mechanism suffices alone---confirming Paper 1's defense-in-depth
architecture.

\textbf{C4: Novel Defense Patterns.} Three domains contributed genuinely
novel extensions to the CIF vocabulary: \emph{verification channel
separation} (Biowarfare), \emph{active perturbation probing} (Trade
Wars), and \emph{physics-informed invariants} (Infrastructure). These
patterns generalize beyond their originating domains and represent
candidate additions to the canonical CIF mechanism set.

\textbf{C5: Temporal Scale Analysis.} The OODA transient dynamics
analysis revealed that Goal Hijacking operates across eight orders of
magnitude in time scale (milliseconds for drone swarms to years for
diplomatic agents), demonstrating that CIF's temporal parameters
(\(\epsilon\), \(\Delta t\)) must be domain-calibrated but the
underlying defense principles are scale-invariant.

\textbf{C6: Real-World Validation.} Retrospective analysis of six
documented AI agent security incidents (2024--2025)---including the
Replit agent meltdown, GitHub Copilot RCE (CVE-2025-53773), Slack AI
data exfiltration, and a \$3.2M procurement fraud---confirms that all
incidents map to one of the three universal attack patterns and would
have been detectable or preventable by the appropriate CIF mechanism.
This provides the first empirical grounding for the CIF-AD-OODA
framework in real production failures (see Supplementary Material S2).
\end{frame}

\begin{frame}[allowframebreaks]{Relationship to the Series}
\protect\phantomsection\label{relationship-to-the-series}
This paper completes a four-part arc in the \emph{Cognitive Security for
Multiagent Operators} series:

\begin{itemize}
\tightlist
\item
  \textbf{Paper 1: Formal Foundations} \cite{friedman2026cogsec1} (DOI:
  10.5281/zenodo.18364119) established the formal foundations: cognitive
  state model
  \(\sigma_i = \langle \mathcal{B}_i, \mathcal{G}_i, \mathcal{I}_i, \mathcal{H}_i \rangle\),
  trust calculus with \(\delta^d\) bounded delegation, adversary
  taxonomy (\(\Omega_1\)--\(\Omega_5\)), information-theoretic
  stealth--impact bounds, and five canonical defense mechanisms with
  composition algebra. A supplementary chapter additionally develops the
  eusocial-colony analogy.
\item
  \textbf{Paper 2: Computational Validation} \cite{friedman2026cogsec2}
  (DOI: 10.5281/zenodo.18364128) provided computational validation:
  benchmark evaluation across 950 attack scenarios, ablation studies,
  Bayesian uncertainty quantification, and colony-scale benchmarks, with
  the recommended defense stack achieving 94--100\% detection in
  parametric simulation.
\item
  \textbf{Paper 3: A Qualitative Review for Practitioners}
  \cite{friedman2026cogsec3} (DOI: 10.5281/zenodo.18364130) translated
  the formal and empirical results into accessible engineering guidance:
  deployment guides, subagent-hardening patterns, incident-response
  playbooks, monitoring strategies, cost--benefit analysis, common
  pitfalls, case studies, and operator risk frameworks. No formal
  prerequisites assumed.
\item
  \textbf{Paper 4: Applications of the Cognitive Integrity Framework}
  (this paper) demonstrates real-world applicability: CIF's formal
  mechanisms map to concrete defenses across ten high-stakes operational
  domains via the integrated CIF-AD-OODA model, and the cross-domain
  analysis yields new insights (three universal attack patterns, three
  novel defense extensions) not visible from any single domain, with
  retrospective validation against six documented 2024--2025 AI agent
  security incidents.
\end{itemize}

Together, the series establishes that cognitive integrity is not merely
a theoretical concern but a \emph{necessary engineering discipline} for
deployed multiagent systems. Readers seeking derivations or proofs
should consult Part 1; readers seeking empirical measurement should
consult Part 2; readers looking to deploy these defenses operationally
should consult Part 3; readers evaluating CIF for a specific operational
sector have just read the present work.
\end{frame}

\begin{frame}[allowframebreaks]{Future Work}
\protect\phantomsection\label{future-work}
Several directions emerge from this analysis:

\begin{enumerate}
\item
  \textbf{Empirical validation.} The most critical next step is
  controlled experimentation in at least one domain---ideally
  cyber-security or infrastructure, where testbed environments
  exist---to validate CIF defense effectiveness against real Goal
  Hijacking attacks with measured detection rates and false positive
  costs. The real-world incidents cataloged in Supplementary Material S2
  provide natural experiment data for retrospective
  validation---particularly the Replit and procurement agent cases,
  where the full attack chain is documented and the hypothesized CIF
  defenses can be simulated against the recorded agent behavior.
\item
  \textbf{Multi-domain attacks.} Adversaries operating across domain
  boundaries (e.g., manipulating food security data to influence trade
  policy agents) represent a class of attacks that single-domain
  analysis cannot capture. Federated CIF architectures with cross-domain
  trust management are needed.
\item
  \textbf{CIF parameter tuning.} Systematic derivation of optimal
  mechanism parameters (\(\tau\), \(\epsilon\), \(q\), \(\Delta t\)) for
  each domain, potentially through automated calibration against
  domain-specific attack distributions.
\item
  \textbf{Automated domain analysis.} The five-step domain analysis
  template is currently applied manually. Automation---where an AI agent
  characterizes its own operational context, identifies its FRs and DPs,
  and selects appropriate CIF mechanisms---would enable self-configuring
  cognitive integrity.
\item
  \textbf{Higher-class adversaries.} Extending the applied analysis to
  \(\Omega_3\)--\(\Omega_5\) attacks and multi-class compositions would
  address the scope limitation identified in
  \cref{sec:limitations_discussion}.
\item
  \textbf{Tool ecosystem security.} The emergence of the Model Context
  Protocol (MCP) and similar agent-tool integration frameworks
  introduces attack surfaces---particularly tool poisoning and tool-call
  interception---not captured by the current \(\Omega_2\) analysis. As
  tool ecosystems become the primary interface between agents and their
  operational environments, CIF mechanisms must be extended to cover the
  tool integration layer explicitly.
\end{enumerate}

\begin{block}{Dual-Use Considerations}
\protect\phantomsection\label{dual-use-considerations}
We note that the CIF-AD-OODA framework, while designed as a defense
methodology, also serves as an analytical tool that could assist
adversaries in identifying undefended attack vectors. The universal
attack pattern taxonomy (C2) and the CIF mechanism coverage matrix
(\cref{sec:mechanism_coverage}) collectively identify which domains are
most vulnerable to which attack patterns and which defense mechanisms
are least deployed. We recommend that practitioners applying this
framework in specific operational domains restrict the detailed defense
mapping to classified or controlled channels, consistent with
responsible disclosure practices in the cybersecurity community.
\end{block}
\end{frame}

\begin{frame}[allowframebreaks]{Closing Statement}
\protect\phantomsection\label{closing-statement}
Prior to this work, the threat of Goal Hijacking was largely viewed as
an issue of content moderation---preventing a chatbot from saying
something inappropriate. We have demonstrated that in the domain of
deployed multiagent systems, Goal Hijacking is a kinetic and existential
threat. It is the ability of an adversary to rewrite the Functional
Requirements of our critical infrastructure, turning our own autonomous
agents against us.

By integrating Axiomatic Design principles with the Cognitive Integrity
Framework and analyzing the temporal dynamics through the OODA lens, we
establish both a theoretical foundation and a practical defense
methodology. We move from fragile ``prompt engineering'' to robust
``goal engineering.'' We secure the OODA loop not by sanitizing the
world of adversarial inputs, but by hardening the agent's Orientation
phase against the transient seduction of the hijack---through Cognitive
Firewalls that filter, Belief Sandboxes that isolate, Behavioral
Invariants that constrain, Drift Detectors that monitor, and Byzantine
Consensus that validates. In doing so, we ensure that as our systems
become faster and more autonomous, they remain unmistakably \emph{ours}.
\end{frame}

\end{document}
